% =========================================================
\section{A completely solvable planar gradient model}
\label{sec:planar-model}
% =========================================================

We now present a two-dimensional model for which the smooth flip
bifurcation, the first clipping contact, the separation coefficient,
and the post-contact dynamics can all be computed explicitly.

Consider the potential
\begin{equation}
V(x,y)
=
\frac{\lambda}{2}x^2
+
\frac{\nu}{2}y^2
-
\frac{b}{4}x^4
+
\frac{d}{6}x^6,
\label{eq:planar-potential}
\end{equation}
where
\begin{equation}
\lambda>\nu>0,
\qquad
b>0,
\qquad
d>0,
\qquad
b^2<4d\lambda.
\label{eq:planar-parameters}
\end{equation}

We use the standard Euclidean clipping norm throughout this section.

The gradient is
\begin{equation}
\nabla V(x,y)
=
\begin{pmatrix}
\lambda x-bx^3+dx^5
\\
\nu y
\end{pmatrix}.
\label{eq:planar-gradient}
\end{equation}
For convenience, define
\begin{equation}
g(x)
:=
\lambda x-bx^3+dx^5
=
x\bigl(\lambda-bx^2+dx^4\bigr).
\label{eq:g-definition}
\end{equation}

The condition $b^2<4d\lambda$ implies
\begin{equation}
\lambda-bs+ds^2>0
\qquad
\text{for every }s\geq0,
\label{eq:g-positive-factor}
\end{equation}
and hence
\begin{equation}
xg(x)>0
\qquad
\text{for every }x\neq0.
\label{eq:g-sign}
\end{equation}

Thus $(0,0)$ is the unique critical point of $V$ and is a strict
global minimum.

% ---------------------------------------------------------
\subsection{Verification of the standing hypotheses}
\label{subsec:planar-hypotheses}
% ---------------------------------------------------------

At the origin,
\begin{equation}
H_*
=
D^2V(0,0)
=
\begin{pmatrix}
\lambda&0
\\
0&\nu
\end{pmatrix}.
\label{eq:planar-Hessian}
\end{equation}
Since
\[
\lambda>\nu>0,
\]
the maximal Hessian eigenvalue is
\begin{equation}
\lambda_*=\lambda
\label{eq:planar-lambda-star}
\end{equation}
and its normalized eigenvector is
\begin{equation}
v_*
=
e_1
=
\begin{pmatrix}
1\\0
\end{pmatrix}.
\label{eq:planar-v-star}
\end{equation}

The flip threshold is therefore
\begin{equation}
\boxed{
\eta_f
=
\frac{2}{\lambda}.
}
\label{eq:planar-flip-threshold}
\end{equation}

Since the potential is even in $x$,
\begin{equation}
D^3V(0,0)=0.
\label{eq:planar-third-zero}
\end{equation}
Moreover,
\begin{equation}
D^4V(0,0)[e_1,e_1,e_1,e_1]
=
-6b.
\label{eq:planar-fourth}
\end{equation}

Hence
\begin{equation}
\Gamma_*
=
-
D^4V(0,0)[e_1^4]
=
6b>0.
\label{eq:planar-Gamma}
\end{equation}
Thus the flip is supercritical.

The first flip coefficient is
\begin{equation}
\boxed{
\ell_*
=
\frac{\Gamma_*}{3\lambda}
=
\frac{2b}{\lambda}.
}
\label{eq:planar-ell}
\end{equation}

The general amplitude formula gives
\begin{equation}
\boxed{
A_*^2
=
\frac{3\lambda^2}{\Gamma_*}
=
\frac{\lambda^2}{2b},
}
\label{eq:planar-A-square}
\end{equation}
or
\begin{equation}
A_*
=
\frac{\lambda}{\sqrt{2b}}.
\label{eq:planar-A}
\end{equation}

Finally, the general Euclidean separation coefficient becomes
\begin{equation}
\boxed{
K_2
=
\frac{\Gamma_*}{3\lambda^4}
=
\frac{2b}{\lambda^4}.
}
\label{eq:planar-K-general-prediction}
\end{equation}

We shall recover all of these quantities below from the exact
dynamics.

% ---------------------------------------------------------
\subsection{Exact smooth gradient dynamics}
\label{subsec:planar-smooth-dynamics}
% ---------------------------------------------------------

The ordinary gradient-descent map is
\begin{equation}
G_\eta(x,y)
=
\begin{pmatrix}
(1-\eta\lambda)x
+
\eta bx^3
-
\eta dx^5
\\[1mm]
(1-\eta\nu)y
\end{pmatrix}.
\label{eq:planar-GD-map}
\end{equation}

The line
\begin{equation}
\mathcal L
=
\{(x,0):x\in\mathbb R\}
\label{eq:invariant-axis}
\end{equation}
is invariant.

On $\mathcal L$, the dynamics reduce to the odd scalar map
\begin{equation}
f_\eta(x)
=
(1-\eta\lambda)x
+
\eta bx^3
-
\eta dx^5.
\label{eq:scalar-smooth-map}
\end{equation}

Because $f_\eta$ is odd, a symmetric period-two orbit has the form
\[
\{-a,a\},
\qquad
a>0.
\]

The condition
\[
f_\eta(a)=-a
\]
is equivalent to
\begin{equation}
\lambda-\frac{2}{\eta}
=
ba^2-da^4.
\label{eq:smooth-cycle-equation}
\end{equation}

Set
\begin{equation}
q:=a^2.
\label{eq:q-amplitude}
\end{equation}
Then
\begin{equation}
dq^2-bq
+
\lambda-\frac{2}{\eta}
=
0.
\label{eq:q-quadratic}
\end{equation}

The small-amplitude root is therefore
\begin{equation}
q_{\mathrm{sm}}(\eta)
=
\frac{
b-
\sqrt{
b^2
-
4d\left(\lambda-\frac{2}{\eta}\right)
}
}{
2d}.
\label{eq:q-smooth-exact}
\end{equation}

Thus the smooth period-two branch is given explicitly by
\begin{equation}
\boxed{
\mathcal O_{2,\mathrm{sm}}(\eta)
=
\left\{
(-a_{\mathrm{sm}}(\eta),0),
(a_{\mathrm{sm}}(\eta),0)
\right\},
}
\label{eq:smooth-cycle-exact}
\end{equation}
where
\begin{equation}
a_{\mathrm{sm}}(\eta)
=
\sqrt{q_{\mathrm{sm}}(\eta)}.
\label{eq:a-smooth}
\end{equation}

% ---------------------------------------------------------
\subsection{Exact verification of the square-root law}
\label{subsec:planar-square-root}
% ---------------------------------------------------------

Let
\begin{equation}
\mu
=
\eta-\frac{2}{\lambda}.
\label{eq:planar-mu}
\end{equation}
Then
\begin{equation}
\lambda-\frac{2}{\eta}
=
\frac{\lambda^2}{2}\mu
-
\frac{\lambda^3}{4}\mu^2
+
O(\mu^3).
\label{eq:lambda-two-over-eta-expansion}
\end{equation}

Solving \eqref{eq:smooth-cycle-equation} gives
\begin{equation}
a_{\mathrm{sm}}(\eta)^2
=
\frac{\lambda^2}{2b}\mu
+
\frac{
\lambda^3(d\lambda-b^2)
}{
4b^3
}
\mu^2
+
O(\mu^3).
\label{eq:a-square-expansion}
\end{equation}

Taking the square root yields
\begin{equation}
a_{\mathrm{sm}}(\eta)
=
\frac{\lambda}{\sqrt{2b}}
\mu^{1/2}
\left[
1
+
\frac{
\lambda(d\lambda-b^2)
}{
4b^2
}
\mu
+
O(\mu^2)
\right].
\label{eq:a-expansion}
\end{equation}

Hence
\begin{equation}
a_{\mathrm{sm}}(\eta)
=
A_*\sqrt{\mu}
+
O(\mu^{3/2}),
\label{eq:planar-square-root-law}
\end{equation}
with
\[
A_*=\frac{\lambda}{\sqrt{2b}},
\]
exactly as predicted by
\eqref{eq:planar-A}.

Because the model is reflection-symmetric, the midpoint of the
smooth two-cycle remains exactly at the critical point:
\begin{equation}
\frac{
(a_{\mathrm{sm}},0)
+
(-a_{\mathrm{sm}},0)
}{2}
=
(0,0).
\label{eq:planar-midpoint-zero}
\end{equation}

% ---------------------------------------------------------
\subsection{Stability of the smooth period-two branch}
\label{subsec:planar-smooth-stability}
% ---------------------------------------------------------

The derivative of the scalar map is
\begin{equation}
f_\eta'(x)
=
1-\eta\lambda
+
3\eta bx^2
-
5\eta dx^4.
\label{eq:scalar-derivative}
\end{equation}

Using the cycle equation
\eqref{eq:smooth-cycle-equation}, one obtains
\begin{equation}
f_\eta'(a)
=
-1
+
2\eta a^2
\left(
b-2da^2
\right).
\label{eq:scalar-derivative-cycle}
\end{equation}

Since $f_\eta'$ is even, the critical multiplier of the two-step
map is
\begin{equation}
\rho_x^{\mathrm{sm}}
=
\left[
-1
+
2\eta a^2
\left(
b-2da^2
\right)
\right]^2.
\label{eq:planar-smooth-x-multiplier}
\end{equation}

For $\mu\to0^+$,
\begin{equation}
\rho_x^{\mathrm{sm}}
=
1-4\lambda\mu
+
O(\mu^2),
\label{eq:planar-smooth-x-expansion}
\end{equation}
in agreement with
Theorem~\ref{thm:period-two-stability}.

The transverse multiplier is
\begin{equation}
\rho_y^{\mathrm{sm}}
=
(1-\eta\nu)^2.
\label{eq:planar-smooth-y-multiplier}
\end{equation}

Since $0<\nu<\lambda$, at $\eta=\eta_f$ one has
$|1-2\nu/\lambda|<1$, so the transverse multiplier stays inside the
unit disk.  Consequently, the smooth period-two cycle is locally
asymptotically stable for all sufficiently small $\eta-\eta_f>0$.

% ---------------------------------------------------------
\subsection{Exact gradient exposure}
\label{subsec:planar-exposure}
% ---------------------------------------------------------

At the positive point of the cycle,
\[
\nabla V(a_{\mathrm{sm}},0)
=
\begin{pmatrix}
g(a_{\mathrm{sm}})
\\
0
\end{pmatrix}.
\]
By the exact period-two relation,
\begin{equation}
g(a_{\mathrm{sm}})
=
\frac{2a_{\mathrm{sm}}}{\eta}.
\label{eq:planar-exact-gradient-balance}
\end{equation}

Hence the common Euclidean gradient exposure is
\begin{equation}
\boxed{
\mathscr E(\eta)
=
\frac{2a_{\mathrm{sm}}(\eta)}{\eta}.
}
\label{eq:planar-exposure-exact}
\end{equation}

Using \eqref{eq:a-expansion},
\begin{equation}
\mathscr E(\mu)
=
\frac{\lambda^2}{\sqrt{2b}}
\mu^{1/2}
\left[
1
+
\frac{
\lambda(d\lambda-3b^2)
}{
4b^2
}
\mu
+
O(\mu^2)
\right].
\label{eq:planar-exposure-expansion}
\end{equation}

In particular,
\begin{equation}
B_*
=
\frac{\lambda^2}{\sqrt{2b}},
\label{eq:planar-B}
\end{equation}
and therefore
\begin{equation}
\frac{1}{B_*^2}
=
\frac{2b}{\lambda^4},
\label{eq:planar-B-inverse}
\end{equation}
again agreeing with the general separation coefficient
\eqref{eq:planar-K-general-prediction}.

% ---------------------------------------------------------
\subsection{Exact switching-contact equation}
\label{subsec:planar-contact-equation}
% ---------------------------------------------------------

The first contact occurs when
\begin{equation}
\frac{2a_{\mathrm{sm}}(\eta)}{\eta}
=
c.
\label{eq:planar-contact-condition}
\end{equation}
Hence at contact
\begin{equation}
a_{\mathrm{sm}}
=
\frac{\eta c}{2}.
\label{eq:planar-contact-amplitude}
\end{equation}

Substituting this into
\eqref{eq:smooth-cycle-equation} gives the exact scalar equation
\begin{equation}
\Psi(\eta,c)
:=
\lambda
-
\frac{2}{\eta}
-
\frac{b}{4}\eta^2c^2
+
\frac{d}{16}\eta^4c^4
=
0.
\label{eq:Psi-contact}
\end{equation}

At
\[
(\eta,c)
=
\left(
\frac{2}{\lambda},0
\right),
\]
we have
\[
\Psi\left(\frac{2}{\lambda},0\right)=0
\]
and
\begin{equation}
\partial_\eta\Psi
\left(
\frac{2}{\lambda},0
\right)
=
\frac{\lambda^2}{2}
>0.
\label{eq:Psi-eta-positive}
\end{equation}

Thus the implicit function theorem gives a unique smooth
first-contact curve
\[
\eta=\eta_{\mathrm{sc}}(c)
\]
for sufficiently small $c$.

% ---------------------------------------------------------
\subsection{Exact verification of the separation law}
\label{subsec:planar-separation}
% ---------------------------------------------------------

Write
\begin{equation}
\eta_{\mathrm{sc}}(c)
=
\frac{2}{\lambda}
+
K_2c^2
+
K_4c^4
+
O(c^6).
\label{eq:planar-eta-series}
\end{equation}

Substituting this expansion into
\eqref{eq:Psi-contact} and comparing powers of $c$ gives
\begin{equation}
\boxed{
K_2
=
\frac{2b}{\lambda^4},
}
\label{eq:planar-K2}
\end{equation}
and
\begin{equation}
\boxed{
K_4
=
\frac{6b^2}{\lambda^7}
-
\frac{2d}{\lambda^6}.
}
\label{eq:planar-K4}
\end{equation}

Therefore
\begin{equation}
\boxed{
\eta_{\mathrm{sc}}(c)
=
\frac{2}{\lambda}
+
\frac{2b}{\lambda^4}c^2
+
\left(
\frac{6b^2}{\lambda^7}
-
\frac{2d}{\lambda^6}
\right)c^4
+
O(c^6).
}
\label{eq:planar-separation-final}
\end{equation}

The leading coefficient
\[
\frac{2b}{\lambda^4}
\]
coincides exactly with
\[
\frac{\Gamma_*}{3\lambda_*^4}
=
\frac{6b}{3\lambda^4},
\]
as predicted by
Theorem~\ref{thm:quadratic-separation}.

This model therefore verifies not only the quadratic exponent but
also the absence of a cubic correction.

% ---------------------------------------------------------
\subsection{The fully clipped symmetric branch}
\label{subsec:planar-clipped-symmetric}
% ---------------------------------------------------------

We now turn to the post-contact regime.  Define
\begin{equation}
a_c(\eta)
:=
\frac{\eta c}{2},
\label{eq:ac-definition}
\end{equation}
and consider the pair
\begin{equation}
Y_\pm(\eta,c)
=
(\pm a_c(\eta),0).
\label{eq:symmetric-clipped-pair}
\end{equation}
Its gradient magnitude is
\begin{equation}
g(a_c)
=
a_c
\left(
\lambda
-
ba_c^2
+
da_c^4
\right),
\label{eq:g-ac}
\end{equation}
and a direct calculation gives
\begin{equation}
g(a_c)-c
=
\frac{\eta c}{2}
\Psi(\eta,c),
\label{eq:g-minus-c-Psi}
\end{equation}
with $\Psi$ as in \eqref{eq:Psi-contact}.  Hence $g(a_c)=c$ at
$\eta=\eta_{\mathrm{sc}}(c)$, and since
$\partial_\eta\Psi(\eta_{\mathrm{sc}}(c),c)>0$ for all sufficiently
small $c$,
\begin{equation}
g(a_c)>c
\qquad
\text{for }
\eta>\eta_{\mathrm{sc}}(c)
\label{eq:ac-fully-clipped}
\end{equation}
sufficiently close to the contact value.  Both points $Y_\pm$
therefore lie in the fully clipped region.  On the invariant axis,
radial clipping acts as
\[
F_{\eta,c}(x,0)
=
(x-\eta c,0)\quad(g(x)>c),
\qquad
F_{\eta,c}(x,0)
=
(x+\eta c,0)\quad(|g(x)|>c,\ x<0),
\]
and since $2a_c=\eta c$,
\begin{equation}
F_{\eta,c}(a_c,0)
=
(-a_c,0),
\qquad
F_{\eta,c}(-a_c,0)
=
(a_c,0).
\label{eq:clipped-symmetric-forward}
\end{equation}

Thus:

\begin{proposition}[Explicit fully clipped period-two branch]
\label{prop:explicit-clipped-branch}
For every sufficiently small $c>0$ and every
$\eta>\eta_{\mathrm{sc}}(c)$ sufficiently close to
$\eta_{\mathrm{sc}}(c)$,
\begin{equation}
\boxed{
\mathcal O_{2,\mathrm{clip}}(\eta,c)
=
\left\{
\left(-\frac{\eta c}{2},0\right),
\left(\frac{\eta c}{2},0\right)
\right\}
}
\label{eq:explicit-clipped-cycle}
\end{equation}
is a fully clipped period-two orbit.

Its amplitude satisfies exactly
\begin{equation}
\boxed{
2a_c=\eta c.
}
\label{eq:exact-amplitude-saturation-planar}
\end{equation}
\end{proposition}

% ---------------------------------------------------------
\subsection{A continuum of clipped period-two cycles}
\label{subsec:continuum-clipped-cycles}
% ---------------------------------------------------------

The symmetric clipped cycle is not isolated.

For
\[
\eta>\eta_{\mathrm{sc}}(c),
\]
we have
\[
g(a_c)>c.
\]
By continuity, there exists $\delta>0$ such that
\begin{equation}
g(a_c+s)>c,
\qquad
g(a_c-s)>c
\label{eq:clipped-neighborhood-condition}
\end{equation}
for every
\[
|s|<\delta.
\]

Define
\begin{align}
Y_+(s)
&=
(a_c+s,0),
\label{eq:Y-plus-s}
\\
Y_-(s)
&=
(-a_c+s,0).
\label{eq:Y-minus-s}
\end{align}

After reducing $\delta$ if necessary,
\[
a_c+s>0,
\qquad
-a_c+s<0.
\]

Since
\[
Y_+(s)-Y_-(s)
=
(\eta c,0),
\]
radial clipping gives
\begin{align}
F_{\eta,c}(Y_+(s))
&=
Y_-(s),
\label{eq:family-forward}
\\
F_{\eta,c}(Y_-(s))
&=
Y_+(s).
\label{eq:family-backward}
\end{align}

We therefore obtain the following stronger result.

\begin{theorem}[A local family of clipped period-two cycles]
\label{thm:family-clipped-cycles}
For every sufficiently small $c>0$ and
$\eta>\eta_{\mathrm{sc}}(c)$ sufficiently close to contact, there
exists $\delta>0$ such that
\begin{equation}
\mathcal O_2(s)
=
\left\{
(-a_c+s,0),
(a_c+s,0)
\right\},
\qquad
|s|<\delta,
\label{eq:period-two-family}
\end{equation}
is a one-parameter family of fully clipped period-two cycles.

Equivalently,
\begin{equation}
F_{\eta,c}^2(x,0)
=
(x,0)
\label{eq:F2-identity-axis}
\end{equation}
on a nontrivial interval of the positive clipped component and on
its image under $F_{\eta,c}$.
\end{theorem}

\begin{proof}
For $|s|<\delta$,
both points lie in the clipped regime by
\eqref{eq:clipped-neighborhood-condition}.

At $Y_+(s)$, the gradient points in the positive $x$ direction, so
the clipped step is exactly
\[
-\eta c e_1.
\]
Therefore
\[
F_{\eta,c}(Y_+(s))
=
(a_c+s-\eta c,0)
=
(-a_c+s,0)
=
Y_-(s).
\]

At $Y_-(s)$, the gradient points in the negative $x$ direction, so
the clipped step is
\[
+\eta c e_1.
\]
Thus
\[
F_{\eta,c}(Y_-(s))
=
(-a_c+s+\eta c,0)
=
(a_c+s,0)
=
Y_+(s).
\]
\end{proof}

\begin{remark}[Geometric meaning of the unit multiplier]
\label{rem:family-unit-meaning}
The continuum in
Theorem~\ref{thm:family-clipped-cycles} gives a concrete nonlinear
interpretation of the universal unit multiplier proved in
Theorem~\ref{thm:universal-unit-main}.

The neutral multiplier is not merely an artifact of linearization:
in this model the two-step map is exactly the identity along the
one-dimensional clipped family.
\end{remark}

% ---------------------------------------------------------
\subsection{Exact post-contact multipliers}
\label{subsec:planar-post-multipliers}
% ---------------------------------------------------------

Consider the period-two cycle
\[
\mathcal O_2(s)
=
\{Y_-(s),Y_+(s)\}.
\]
Set
\begin{equation}
r_+(s)
=
g(a_c+s),
\qquad
r_-(s)
=
g(a_c-s).
\label{eq:r-plus-minus}
\end{equation}

At either point, the Euclidean clipping direction is parallel to
$e_1$. Therefore
\begin{equation}
I-e_1e_1^\top
=
\begin{pmatrix}
0&0\\
0&1
\end{pmatrix}.
\label{eq:e1-projector}
\end{equation}

Since
\[
D^2V(x,0)
=
\begin{pmatrix}
\lambda-3bx^2+5dx^4&0
\\
0&\nu
\end{pmatrix},
\]
the one-step Jacobians along the clipped cycle are
\begin{align}
DF_{\eta,c}(Y_+(s))
&=
\begin{pmatrix}
1&0
\\
0&
1-\dfrac{\eta\nu c}{r_+(s)}
\end{pmatrix},
\label{eq:Jac-plus-family}
\\
DF_{\eta,c}(Y_-(s))
&=
\begin{pmatrix}
1&0
\\
0&
1-\dfrac{\eta\nu c}{r_-(s)}
\end{pmatrix}.
\label{eq:Jac-minus-family}
\end{align}

Hence the two-step monodromy matrix is
\begin{equation}
M_{\mathrm{clip}}(s)
=
\begin{pmatrix}
1&0
\\
0&
\displaystyle
\left(
1-\frac{\eta\nu c}{r_+(s)}
\right)
\left(
1-\frac{\eta\nu c}{r_-(s)}
\right)
\end{pmatrix}.
\label{eq:family-monodromy}
\end{equation}

Thus the two multipliers are
\begin{equation}
\boxed{
\rho_x^{\mathrm{clip}}=1
}
\label{eq:planar-unit-multiplier}
\end{equation}
and
\begin{equation}
\boxed{
\rho_y^{\mathrm{clip}}(s)
=
\left(
1-\frac{\eta\nu c}{r_+(s)}
\right)
\left(
1-\frac{\eta\nu c}{r_-(s)}
\right).
}
\label{eq:planar-transverse-multiplier}
\end{equation}

For the symmetric member $s=0$,
\begin{equation}
\rho_y^{\mathrm{clip}}(0)
=
\left(
1-\frac{\eta\nu c}{g(a_c)}
\right)^2.
\label{eq:symmetric-clipped-transverse}
\end{equation}

At the switching contact,
\[
g(a_c)=c,
\]
so the outer limiting transverse multiplier becomes
\begin{equation}
\rho_y^{\mathrm{clip}}
=
(1-\eta_{\mathrm{sc}}\nu)^2.
\label{eq:transverse-contact-limit}
\end{equation}

Since
\[
\eta_{\mathrm{sc}}(c)
\to
\frac{2}{\lambda}
\]
as $c\to0^+$ and
\[
0<\nu<\lambda,
\]
one has
\[
\left|
1-\eta_{\mathrm{sc}}(c)\nu
\right|
<1
\]
for all sufficiently small $c$.

Hence the clipped family is transversely attracting while remaining
exactly neutral along its one-dimensional period-two direction.

% ---------------------------------------------------------
\subsection{The local transition in the solvable model}
\label{subsec:planar-transition-summary}
% ---------------------------------------------------------

The planar model provides an explicit realization of the entire
smooth-to-clipped scenario.  For $\eta<\eta_f=2/\lambda$ the origin
is locally asymptotically stable; for
$\eta_f<\eta<\eta_{\mathrm{sc}}(c)$ there is a unique small stable
smooth two-cycle
$\{(-a_{\mathrm{sm}},0),(a_{\mathrm{sm}},0)\}$ whose amplitude
grows as
\[
a_{\mathrm{sm}}
\sim
\frac{\lambda}{\sqrt{2b}}
\sqrt{\eta-\eta_f}.
\]
At $\eta=\eta_{\mathrm{sc}}(c)$ both points reach the clipping
boundary simultaneously with
\[
a_{\mathrm{sm}}
=
\frac{\eta_{\mathrm{sc}}c}{2},
\qquad
\eta_{\mathrm{sc}}(c)
=
\frac{2}{\lambda}
+
\frac{2b}{\lambda^4}c^2
+
O(c^4).
\label{eq:planar-summary-separation}
\]
For $\eta>\eta_{\mathrm{sc}}(c)$ sufficiently close to contact, the
clipped map possesses a one-parameter family of period-two cycles
with exact step constraint $x_+-x_-=\eta c$, and the smooth critical
multiplier $\rho_x^{\mathrm{sm}}<1$ is replaced by the exact neutral
multiplier $\rho_x^{\mathrm{clip}}=1$.  The model thus exhibits the
transition from a stable smooth period-two branch, through a
simultaneous clipping contact, to a neutral family of clipped
two-cycles.

\subsection{A concrete parameter choice}

To make the preceding formulas explicit, consider the parameter
choice
\[
\lambda=2,\qquad
\nu=\frac12,\qquad
b=1,\qquad
d=1.
\]
Then
\[
\eta_f=1,
\qquad
\Gamma_*=6,
\qquad
\ell_*=1,
\]
and therefore
\[
A_*=\sqrt{2},
\qquad
B_*=2\sqrt{2},
\qquad
K_2=\frac18.
\]
Hence the smooth period-two branch satisfies
$x_\pm(\mu)=x_*\pm\sqrt{2\mu}\,v_*+O(\mu)$, while the first
switching threshold obeys
\[
\eta_{\mathrm{sc}}(c)
=
1+\frac18c^2+O(c^4),
\]
with the next correction computable explicitly:
\[
\eta_{\mathrm{sc}}(c)
=
1+\frac18c^2+\frac1{64}c^4+O(c^6).
\]
The example thus provides an explicit realization of the general
quadratic separation law of Section~\ref{sec:separation-law}.